\documentclass[10pt]{article}
\usepackage[margin=24mm]{geometry}
\usepackage{amsmath,booktabs,hyperref}
\title{Re-executing Scientific Design after Method Changes\\\large Public development manuscript skeleton, 19 October 2026}
\author{Working draft}
\date{}
\begin{document}
\maketitle
\begin{abstract}
Changing a device, changing the method that designs it, and changing the procedure
that improves that method invalidate different computations. We study these three
interventions separately in an executable battery-free design apparatus. This
development draft reports six deterministic invalidation controls and four fresh
prompt-intervention attempts. Three delivered qualified computational designs; the
fourth stopped at a frozen resource gate before validation. All 68 model requests
settled. The method update remains unpromoted. No method advantage, recursive
self-improvement, unseen-task transfer, hardware validation, or speedup is established.
\end{abstract}

\section{Question and scope}
A previously successful device trace can cease to describe the system after a
legal change. Reusing that trace may silently preserve obsolete control decisions.
At the method level, a changed global design prompt may alter the first observation,
the candidate, and every later experiment. Evaluating the old device under that
prompt does not evaluate the changed method. Our question is which results remain
valid after a typed intervention, and how to confirm a proposed method change
without charging a cached response as a new design attempt.

\section{Objects and validity contract}
Let $D$ contain device structure, circuits, firmware, storage, communication, and
deployment. Let $H$ contain the design prompt, skills, available observations,
context policy, experiment order, and organization. Let $G$ propose and select
changes to $H$. Device runtime, the agent's decision span, and method generations
are separate quantities. A record is reusable only when all consumed inputs,
environment identities, numerical settings, and output hashes remain applicable.
An evaluation change may reuse raw physical output while requiring a new verdict.

For a global prompt change $H\to H'$, our conservative rule restarts at task entry
with a new owner, initial candidate, and context. Exact lower-level input responses
may still be reused. A local continuation requires a separate proof that its
restored state is valid. A $G$ comparison additionally fixes the starting $H_0$,
development distribution, selection protocol, and improvement budget.

\section{F1a: device-response controls}
We compare old fixed-trajectory replay, conservative dependency invalidation,
boundary invalidation, full response execution, and exact-input response caching.
The first qualification apparatus has energy arrivals, finite storage, a threshold
controller, and load consumption. The independent full response is the oracle.
In a feedback case, both dependency and boundary invalidation expand to the full
response; this null advantage is a required control, not a speedup result.

In the executed deterministic controls, fixed-trajectory replay disagrees with the full response on 5 of 6 inputs. The conservative dependency, boundary, and exact-cache results agree on all 6. Only the unchanged input is an initial exact-cache hit. These are synthetic correctness controls; execution counts do not measure wall-time speedup.
\begin{center}
\begin{tabular}{lrrr}
\toprule
Input & Replay packets & Full packets & Replay valid\\
\midrule
unchanged & 4 & 4 & True\\
less-energy & 4 & 2 & False\\
higher-load & 4 & 3 & False\\
threshold-only & 4 & 3 & True\\
capacity-limit & 4 & 0 & False\\
arrival-order & 4 & 2 & False\\
\bottomrule
\end{tabular}
\end{center}


\section{F1b and early F4: real method re-execution}
The frozen public task uses the already qualified piezoelectric, vendor PMIC,
host-C controller, and digital receiver apparatus. The power board is fixed; the
agent must revise firmware and deliver six valid frames in twenty seconds.
The declared initial charge and current envelope remain assumptions. The task is
an exposed development derivative of historical T1, not a fresh holdout.

$H_0$ requests one complete design attempt with execution, validation, decision,
report, release, and close. $H_1$ adds a constraint and time-origin audit before a
single justified revision, with version-bound evidence. Both receive the same
full task and sources, backend, tool definitions, limits, and initial native cache.
The prompt update was proposed by the research operator from earlier timing and
tool-use failures; it is not an autonomous improver. Two attempts per condition run
in order $H_0,H_1,H_1,H_0$. Each allows at most 80 parent-plus-child requests,
four candidate evaluations, and 1,200 seconds of decision wall time. Organization
is single-agent. No automatic resampling is allowed.

The protocol requires full execution to confirm an attempted method. Development retention requires
both $H_1$ confirmations to deliver admissible results without an increased
failure count. With two attempts per condition, cost and quality differences are
descriptive. A tie, failure, or no retained update is a valid result.

\paragraph{Observed attempts and resource stop.} All 4 fresh controllers consumed their assigned goal, started with zero prior trajectory messages, and changed the initial firmware. After administrative identifiers are normalized, the base prompts match; the harness, task, initial device, and physical cache start are identical. 3 attempts completed all six computational qualification gates and delivery. The last attempt, $H_0$-r2, completed its native numerical execution but hit the frozen 5 GB free-space gate before validation and release. It stopped after 260.958 seconds and was not resampled. Its pending qualification is not converted into a device failure or a pass.
\begin{center}
{\small\begin{tabular}{lrrrrrl}
\toprule
Attempt & Requests & USD upper & Executions & Native & Frames & Outcome\\
\midrule
H0-r1 & 13 & 0.188330 & 1 & 2 & 14 & PASS\\
H1-r1 & 20 & 0.358763 & 2 & 2 & 10 & PASS\\
H1-r2 & 20 & 0.361660 & 2 & 2 & 7 & PASS\\
H0-r2 & 15 & 0.234168 & 2 & 2 & -- & STOP\\
\bottomrule
\end{tabular}}
\end{center}
Across all attempted owners, including the stop, 68 requests settled with no missing provider usage or unresolved reservation. Their conservative charge upper bound is USD 1.142922; the rate-based cache-aware estimate is USD 0.350731, not an invoice. The 7 candidate executions used 8 new native dispatches. Each newly solved input is disjoint from earlier successful native inputs; applicable exact inputs were reused. The sixteen-entry initial cache was copied without new solves. Its historical construction bill is unknown, not zero. Operator proposal and Codex supervision, executor-chat, and writing costs are outside the BF ledger; these attempt charges cannot establish end-to-end improvement efficiency.

Both $H_1$ confirmations delivered, using twenty requests each; the completed $H_0$-r1 used thirteen. The interrupted second $H_0$ prevents a completed arm comparison. We preserve $H_1$ as an experimental artifact, keep the default method unchanged, and make no success-rate, cost, speedup, transfer, or recursive improvement claim. Provider and OS caches were uncontrolled; setup and subsequent wall times remain separate in the audit.


\section{Costs, exposure, and competing explanations}
Report all parent and child requests, settled usage, reservations, failed tools,
native dispatches, retries, confirmations, cache construction, and unresolved
UNKNOWN separately. Total cost is
\[
C=C_{\rm proposal}+C_{\rm selection}+C_{\rm confirmation}
  +C_{\rm failures}+C_{\rm cache\ construction}.
\]
A valid cache hit does not remove the cost of building its entry. Paired attempts
start from identical cache manifests rather than warming later conditions with
earlier outputs. Provider identity and usage availability are dispatch gates.
Preflight failures remain recorded, and cannot become successful attempts by
silent replacement. Mechanism exposure is shared with the interface study.

F1a supports a correctness boundary for this small control; it does not establish
a general incremental algorithm. F1b can establish that a prompt was actually
changed and new design decisions executed. Favorable development outcomes would
still admit task familiarity, stochastic variation, and operator selection as
competing explanations. Later selection and final tasks must be frozen separately.

\section{Related work and remaining experiments}
Executable build-system decompositions provide a strong dependency-invalidation
baseline~\cite{mokhov2018}. Prompt optimizers such as GEPA propose and evaluate
prompt changes using execution feedback~\cite{agrawal2026gepa}. Our current work
does not establish superiority over either family. The planned study isolates
physical-response reuse from full method confirmation and retains negative cases.
F2/F3 will explain decision spans and organizations on a useful subset. F5 requires
independent $G_0/G_1$ lineages from a fixed $H_0$; no such lineage experiment is
reported. F6 will follow the continuous A/R product after its service reference is
qualified. Physical calibration and hardware acceptance remain separate work.

\begin{thebibliography}{9}
\bibitem{mokhov2018} Andrey Mokhov, Neil Mitchell, and Simon Peyton Jones.
Build systems a la carte. ICFP, 2018.
\url{https://www.microsoft.com/en-us/research/publication/build-systems-la-carte/}.
\bibitem{agrawal2026gepa} Lakshya A. Agrawal et al.
GEPA: Reflective Prompt Evolution Can Outperform Reinforcement Learning.
arXiv:2507.19457v2, 2026.
\url{https://arxiv.org/abs/2507.19457v2}.
\end{thebibliography}
\end{document}
